The geometry of the target is either described in terms of regions or in terms 
of cells. In 1-D, the geometry of regions is simply defined by the positions of 
the surfaces and the interfaces between the regions (\texttt{POSIF}), e.g.,
%
\begin{verbatim}
 &geometry posif=-10,0,10000 /
\end{verbatim}
%
Cells are usually finer-grained than regions. Only regions are specified on the
\texttt{\&GEOMETRY} record. Cells are either defined by reading a \texttt{CELL}
file, or they are internally generated from regions in case of a dynamic
simulation (\texttt{LDYN=T}).

In 2-D, the geometry of each region can be specified either by a closed polygon 
or a signed distance function (SDF). To define the polygons, it is first 
necessary to specify the coordinates of the vertices (corner points) using the 
\texttt{POINT} index variable, e.g.,
%
\begin{verbatim}
 &geometry point=1 pos=0,0 /
 &geometry point=2 pos=-5000,0 /
 ...
\end{verbatim}
%
and then to specify the polygons by listing the point indices in the desired
sequence, using the \texttt{REGION} index variable:
%
\begin{verbatim}
 &geometry region=0 points=1,2,3,4,5,6,7,8 /
 &geometry region=1 points=1,4,5,6 /
 ...
\end{verbatim}
%
Note that the circumference of the complete simulation area can be specified as
region 0. If it is not specified, it is determined automatically from the
other regions. There must not be undefined areas within the circumference, and 
no regions may overlap.

Alternatively, the geometry of each region as well as of the union of all
regions can be defined by SDFs. SDFs must be calculated externally, and are
provided as files, see the parameter \texttt{SDFFILE}.

In both cases of 2-D geometry definition, IMSIL uses an internal mesh 
containing the minimum distance of each cell from the boundary. It is used to 
reduce the number of checks that are necessary to determine whether the ion 
leaves a region. In case of geometry definition by SDFs, the internal mesh
coincides with the mesh of the SDFs.

3-D geometry definition is possible via SDFs only.

It is possible to impose cylindrical (\texttt{LCYLX=T}, \texttt{LCYLY=T}, or
\texttt{LCYLZ=T}) or spherical (\texttt{LSPHER=T}) symmetry. Specification of
cylindrical geometry requires 1-D or 2-D geometry definition, where the 
distance along the axis of symmetry replaces the $x$ coordinate,
while the distance from the axis replaces the $z$ coordinate. Specification of
spherical symmetry can only be combined with 1-D geometry definition, where 
the distance from the center of symmetry replaces the $z$ coordinate. 

\begin{keydescription}{\texttt{BNDTYPEX} --- x boundary type}
%
  This parameter specifies the type of boundary at the minimum and maximum $x$
  coordinates of the geometry, where \texttt{BNDTYPEX(1)} refers to the minimum
  $x$ coordinate and \texttt{BNDTYPEX(2)} to the maximum $x$ coordinate.
  Possible values are \texttt{'free'} (free surface through which projectiles
  may enter and leave the target), \texttt{'periodic'} (the geometry is 
  repeated), \texttt{'reflective'} (the geometry is reflected), and 
  \texttt{'extending'} (the geometry as it meets the boundary is extended to 
  plus/minus infinity perpendicular to the boundary). In the combination
  \texttt{'periodic'-'reflective'} the geometry is first reflected 
  and the resulting geometry is repeated. The combinations 
  \texttt{'periodic'-'free'}, \texttt{'periodic'-'extending'}, 
  and \texttt{'reflective'-'reflective'} are forbidden. \texttt{BNDTYPEX} 
  is ignored unless $\texttt{NDIM} \ge 2$.  
  \begin{keytab}
    Type:    \> simple array (2) of character(80) \\
    Default: \> \texttt{'free'}, \texttt{BNDTYPEX(1)} \\
    Range:   \> \texttt{'free'}, \texttt{'periodic'}, \texttt{'reflective'}, 
                \texttt{'extending'} 
  \end{keytab}
\end{keydescription}

\begin{keydescription}{\texttt{BNDTYPEY} --- y boundary type}
%
  This parameter specifies the type of boundary at the minimum and maximum $y$
  coordinates of the geometry, where \texttt{BNDTYPEY(1)} refers to the minimum
  $y$ coordinate and \texttt{BNDTYPEY(2)} to the maximum $y$ coordinate.
  Possible values are \texttt{'free'} (free surface through which projectiles
  may enter and leave the target), \texttt{'periodic'} (the geometry is 
  repeated), \texttt{'reflective'} (the geometry is reflected), and 
  \texttt{'extending'} (the geometry as it meets the boundary is extended to 
  plus/minus infinity perpendicular to the boundary). In the combination
  \texttt{'periodic'-'reflective'} the geometry is first reflected 
  and the resulting geometry is repeated. The combinations 
  \texttt{'periodic'-'free'}, \texttt{'periodic'-'extending'}, 
  and \texttt{'reflective'-'reflective'} are forbidden. \texttt{BNDTYPEY} 
  is ignored unless $\texttt{NDIM} = 3$.  
  \begin{keytab}
    Type:    \> simple array (2) of character(80) \\
    Default: \> \texttt{'free'}, \texttt{BNDTYPEY(1)} \\
    Range:   \> \texttt{'free'}, \texttt{'periodic'}, \texttt{'reflective'}, 
                \texttt{'extending'} 
  \end{keytab}
\end{keydescription}

\begin{keydescription}{\texttt{BNDTYPEZ} --- z boundary type}
%
  This parameter specifies the type of boundary at the minimum and maximum $z$
  coordinates of the geometry, where \texttt{BNDTYPEZ(1)} refers to the minimum
  $z$ coordinate and \texttt{BNDTYPEZ(2)} to the maximum $z$ coordinate.
  Possible values are \texttt{'free'} (free surface through which projectiles
  may enter and leave the target), \texttt{'periodic'} (the geometry is 
  repeated), \texttt{'reflective'} (the geometry is reflected), and 
  \texttt{'extending'} (the geometry as it meets the boundary is extended to 
  plus/minus infinity perpendicular to the boundary). In the combination
  \texttt{'periodic'-'reflective'} the geometry is first reflected 
  and the resulting geometry is repeated. The combinations 
  \texttt{'periodic'-'free'}, \texttt{'periodic'-'extending'}, 
  and \texttt{'reflective'-'reflective'} are forbidden. \texttt{BNDTYPEZ(1)}
  should normally be \texttt{'free'} if an external beam is used. 
  \begin{keytab}
    Type:    \> simple array (2) of character(80) \\
    Default: \> \texttt{'free'}, \texttt{BNDTYPEZ(1)} \\
    Range:   \> \texttt{'free'}, \texttt{'periodic'}, \texttt{'reflective'}, 
                \texttt{'extending'} 
  \end{keytab}
\end{keydescription}


\begin{keydescription}{\texttt{CENTER} --- Center of cylindrical and spherical 
regions}
%
  This parameter specifies the center of the cylindrical or spherical regions,
  if \texttt{LCYLX=T} or \texttt{LCYLY=T} or \texttt{LCYLZ=T} or 
  \texttt{LSPHER=T}. The three components of \texttt{CENTER} correspond to $x$, 
  $y$, and $z$ coordinate of the center. For cylindrical symmetries, the
  symmetry line passes through the position specified by \texttt{CENTER} and is
  parallel to the $x$, $y$, or $z$ axis depending on whether the \texttt{LCYLX},
  \texttt{LCYLY}, or \texttt{LCYLZ} flag is set.
  \begin{keytab}
    Type:    \> simple array (3) of real \\
    Default: \> 0., 0., 0. \\
    Range:   \> arbitrary
  \end{keytab}
\end{keydescription}

\begin{keydescription}{\texttt{FRACUPD} --- Minimum fractional density
change for cell update}
%
  A mesh update is triggered in a 1-D dynamic simulation if the change of the
  cell width relative to the nominal cell width (\texttt{WBOX} of the
  \texttt{\&OUTPUT} record) would exceed \texttt{FRACUPD} in any cell. The cell
  widths change due to relaxation of density changes. The condition is only
  checked every \texttt{NIONUPD} ions. 
  \begin{keytab}
    Type:    \> integer \\
    Default: \> 0.1 \\
    Range:   \> $\ge 0$
  \end{keytab}
\end{keydescription}

\ifprivate
\begin{keydescription}{\texttt{FREEX} --- Free boundaries in x-direction}
%
  This parameter specifies whether the left/right boundary of the
  simulation domain is free to move during relaxation steps. The simulation
  domain is the smallest possible rectangle, consistent with the grid spacings,
  containing the target. Therefore, this boundary condition only applies to
  target boundaries with (close to) the smallest/largest coordinates.
  \begin{keytab}
    Type:    \> simple array (2) of logical \\
    Default: \> \texttt{(F,F)} \\
    Range:   \> \texttt{T}, \texttt{F} 
  \end{keytab}
\end{keydescription}

\begin{keydescription}{\texttt{FREEZ} --- Free boundaries in z-direction}
%
  This parameter specifies whether the top/bottom boundary of the
  simulation domain is free to move during relaxation steps. The simulation
  domain is the smallest possible rectangle, consistent with the grid spacings,
  containing the target. Therefore, this boundary condition only applies to
  target boundaries with (close to) the smallest/largest coordinates.
  \begin{keytab}
    Type:    \> simple array (2) of logical \\
    Default: \> \texttt{(T,F)} \\
    Range:   \> \texttt{T}, \texttt{F} 
  \end{keytab}
\end{keydescription}
\fi

\begin{keydescription}{\texttt{GEOMFILE} --- Geometry file}
%
  This parameter specifies the name of a geometry file, which contains further 
  \texttt{\&GEOMETRY} records with \texttt{POINT} and \texttt{REGION}
  definitions. This has the same effect as if the contents of the geometry file
  was inserted before the first \texttt{\&GEOMETRY} record of the input file.
  As a consequence, any specifications in the input file will overwrite
  corresponding specifications in the geometry file. Any other content of the
  geometry file other than \texttt{\&GEOMETRY} records is ignored. If multiple
  geometry files are specified, only the last one will be included. The geometry
  file is intended for the specification of 2-D geometries (points and regions),
  so one geometry file can be used for multiple input files. If a geometry file
  is specified, the geometry points and regions are not written to the output
  file.
  \texttt{GEOMFILE} is ignored for 1-D geometries (\texttt{NDIM=1}).
  \begin{keytab}
    Type:    \> character string \\
    Default: \> -- \\
    Range:   \> Any valid file name ($\le 80$ characters)
  \end{keytab}
\end{keydescription}

\begin{keydescription}{\texttt{GRID} --- Grid type}
%
  This parameter specifies the type of the grid defining the cells for a
  dynamic simulation (\texttt{LDYN=T}). \texttt{GRID='Cartesian'} can be used
  for \texttt{GRID='regular'} for backward compatibility. \texttt{GRID} is
  ignored for \texttt{LDYN=F}.
  \begin{keytab}
    Type:    \> character(80) \\
    Default: \> \texttt{'regular'} \\
    Range:   \> \texttt{'regular'}, \texttt{'rectilinear'} for \texttt{NDIM=1} \\ 
             \> \texttt{'regular'}, \texttt{'curvilinear'} for \texttt{NDIM=2}
  \end{keytab}
\end{keydescription}

\begin{keydescription}{\texttt{LCYLX} --- Cylinder region flag, axis parallel $x$}
%
  This parameter specifies that the target should be considered cylindrical 
  symmetric. The axis of rotation is parallel to the $x$ axis and goes through 
  \texttt{CENTER}. For \texttt{NDIM=1} the radii of the 
  cylindrical regions are specified by \texttt{POSIF}. For \texttt{NDIM=2}
  the polygon defined by \texttt{POS} and \texttt{POINTS} is rotated around 
  the axis. \texttt{POS(1)} is considered the coordinate along the axis, 
  \texttt{POS(2)} the distance from the axis (radius). Any combination of two
  or more of \texttt{LCYLX=T}, \texttt{LCYLY=T}, \texttt{LCYLZ=T}, 
  \texttt{LSPHER=T}, and \texttt{LROBIN=T} is illegal. \texttt{LCYLX=T} can 
  only be specified for \texttt{LDYN=F} and \texttt{NDIM=1} or \texttt{NDIM=2}. 
  \begin{keytab}
    Type:    \> logical \\
    Default: \> \texttt{F} \\
    Range:   \> \texttt{T}, \texttt{F} 
  \end{keytab}
\end{keydescription}

\begin{keydescription}{\texttt{LCYLY} --- Cylinder region flag, axis parallel $y$}
%
  This parameter specifies that the target should be considered cylindrical 
  symmetric. The axis of rotation is parallel to the $x$ axis and goes through 
  \texttt{CENTER}. For \texttt{NDIM=1} the radii of the 
  cylindrical regions are specified by \texttt{POSIF}. For \texttt{NDIM=2}
  the polygon defined by \texttt{POS} and \texttt{POINTS} is rotated around 
  the axis. \texttt{POS(1)} is considered the coordinate along the axis, 
  \texttt{POS(2)} the distance from the axis (radius). Any combination of two
  or more of \texttt{LCYLX=T}, \texttt{LCYLY=T}, \texttt{LCYLZ=T}, 
  \texttt{LSPHER=T}, and \texttt{LROBIN=T} is illegal. \texttt{LCYLY=T} can 
  only be specified for \texttt{LDYN=F} and \texttt{NDIM=1} or \texttt{NDIM=2}.
  \begin{keytab}
    Type:    \> logical \\
    Default: \> \texttt{F} \\
    Range:   \> \texttt{T}, \texttt{F} 
  \end{keytab}
\end{keydescription}

\begin{keydescription}{\texttt{LCYLZ} --- Cylinder region flag, axis parallel $z$}
%
  This parameter specifies that the target should be considered cylindrical 
  symmetric. The axis of rotation is parallel to the $x$ axis and goes through 
  \texttt{CENTER}. For \texttt{NDIM=1} the radii of the 
  cylindrical regions are specified by \texttt{POSIF}. For \texttt{NDIM=2}
  the polygon defined by \texttt{POS} and \texttt{POINTS} is rotated around 
  the axis. \texttt{POS(1)} is considered the coordinate along the axis, 
  \texttt{POS(2)} the distance from the axis (radius). Any combination of two
  or more of \texttt{LCYLX=T}, \texttt{LCYLY=T}, \texttt{LCYLZ=T}, 
  \texttt{LSPHER=T}, and \texttt{LROBIN=T} is illegal. \texttt{LCYLZ=T} can 
  only be specified for \texttt{LDYN=F} and \texttt{NDIM=1} or \texttt{NDIM=2}.
  \begin{keytab}
    Type:    \> logical \\
    Default: \> \texttt{F} \\
    Range:   \> \texttt{T}, \texttt{F} 
  \end{keytab}
\end{keydescription}

\begin{keydescription}{\texttt{LSPHER} --- Spherical region flag}
%
  This parameter specifies that the target regions should be
  considered spherical around the center given by \texttt{CENTER} rather than 
  planar. The radii of the spherical regions are specified by \texttt{POSIF}. 
  Any combinations of \texttt{LCYLX=T}, \texttt{LCYLY=T}, \texttt{LCYLZ=T}, 
  \texttt{LSPHER=T}, and \texttt{LROBIN=T} must not be specified. 
  \texttt{LSPHER=T} can only be specified for \texttt{NDIM=1}, \texttt{LDYN=F}.
  \begin{keytab}
    Type:    \> logical \\
    Default: \> \texttt{F} \\
    Range:   \> \texttt{T}, \texttt{F} 
  \end{keytab}
\end{keydescription}

\begin{keydescription}{\texttt{NBOX} --- Maximum number of boxes}
%
  This parameter specifies the maximum number of boxes (cells) of the 1-D cell 
  grid used to store densities in a dynamic simulation (\texttt{LDYN=T}), if 
  the geometry is defined by regions (\texttt{USECELL=F}). 
  Note that an excessive number of cells increases noise in densities.
  \texttt{NBOX} is effective only if the number of boxes necessary to cover 
  the geometry with a box width of \texttt{WBOX} would exceed \texttt{NBOX}. 
  In this case the box width is adjusted so that the number of boxes equals 
  \texttt{NBOX}. 
  Thus, \texttt{NBOX} sets the upper limit to the number of boxes used in the 
  simulation. 
  \texttt{NBOX=0} means that the number of boxes is determined by 
  \texttt{WBOX} alone. 
  If reflective boundaries are specified, \texttt{NBOX} refers to the number
  of boxes in the unreflected geometry, i.e., it is effectively doubled.
  \texttt{NBOX} is ignored if the geometry is read from a cell file 
  (\texttt{USECELL=T}), or if $\texttt{NDIM} \ne 1$.
  \begin{keytab}
    Type:    \> integer \\
    Default: \> \texttt{NBOX} of \texttt{\&OUTPUT} for \texttt{LDYN=T} if 
                defined \\
             \> 32 otherwise \\
    Range:   \> $\ge 0$
  \end{keytab}
\end{keydescription}


\begin{keydescription}{\texttt{NBOX2} --- Maximum number of boxes}
%
  This parameter specifies the maximum number of boxes of the 2-D cell grid 
  used to store densities in a dynamic simulation (\texttt{LDYN=T}), if the 
  geometry is defined by regions (\texttt{USECELL=F}). 
  It is also used to specify the number of boxes of the internal mesh for both
  static (\texttt{LDYN=F}) and dynamic (\texttt{LDYN=T}) simulations if the
  geometry is defined by regions.
  The internal mesh is used to store the signed distance function that
  internally represents the geometry.
  A finer mesh means a finer resolution of the geometry.
  An excessive number of mesh boxes, however, slows down the simulation, and
  increases noise in a dynamic simulation.
  The first value of \texttt{NBOX2} corresponds to the number of boxes in 
  $x$ direction, the second value to the number of boxes in $z$ direction. 
  Each of the values of \texttt{NBOX2} is effective only if the number of 
  boxes necessary to cover the geometry with a box width of the respective 
  value of \texttt{WBOX2} would exceed \texttt{NBOX2}. 
  In this case the box width is adjusted so that the number of boxes equals 
  \texttt{NBOX2}. 
  Thus, \texttt{NBOX2} sets upper limits to the number of boxes used in the 
  simulation. 
  \texttt{NBOX2=0} means that the number of boxes is determined by 
  \texttt{WBOX2} alone. 
  If reflective boundaries are specified, \texttt{NBOX2} refers to the number
  of boxes in the unreflected geometry, i.e., it is effectively doubled in 
  the respective direction.
  \texttt{NBOX2} is ignored if the geometry is read from an SDF file 
  (\texttt{SDFFILE}) or from a cell file (\texttt{USECELL=T}), or if 
  $\texttt{NDIM} \ne 2$.
  \begin{keytab}
    Type:    \> simple array (2) of integer \\
    Default: \> \texttt{NBOX2} of \texttt{\&OUTPUT} for \texttt{LDYN=T} if
                defined \\
             \> 32, \texttt{NBOX2(1)} for \texttt{LDYN=T} otherwise \\
             \> 1000, \texttt{NBOX2(1)} for \texttt{LDYN=F} \\
    Range:   \> $\ge 0$
  \end{keytab}
\end{keydescription}

\begin{keydescription}{\texttt{POINT} --- Point index}
%
  This parameter specifies the index of a point whose coordinates are specified
  on the same record by \texttt{POS}. The point indices are referred to by the
  \texttt{POINTS} arrays. 
  \begin{keytab}
    Type:    \> integer \\
    Default: \> -- \\
    Range:   \> $>0$ 
  \end{keytab}
\end{keydescription}

\begin{keydescription}{\texttt{POINTS} --- Points defining a region polygon}
%
  This parameter specifies the indices of the points (as defined by
  \texttt{POINT}) defining the region polygon.  Specification of
  \texttt{POINTS} requires the index variable \texttt{REGION} on the same
  record. A maximum of 20 points may be specified on one record.
  Specifications of \texttt{POINTS} on different records for the same
  \texttt{REGION} are joined together.  In this way, larger arrays of points 
  may be split to several records.  
  \begin{keytab}
    Type:    \> index variable array (\texttt{REGION}) of arbitrary-size 
    simple arrays of integer \\
    Default: \> -- \\
    Range:   \> point indices defined by \texttt{POINT} 
  \end{keytab}
\end{keydescription}

\begin{keydescription}{\texttt{POS} --- Position of an interface point}
%
  This parameter specifies the position of an interface point. The
  position is a simple array of size 2, consisting of the $x$ and $z$
  coordinate in case of Cartesian coordinates. In case of cylindrical
  coordinates \texttt{POS(1)} refers to the distance along the axis of
  symmetry, and \texttt{POS(2)} to the distance from the axis of symmetry.
  The units of the coordinates are \AA.  Specification of \texttt{POS} 
  requires the index variable \texttt{POINT} on the same record. 
  \begin{keytab}
    Type:    \> index variable array (\texttt{POINT}) of simple array
                (2) of real \\
    Default: \> -- \\
   Range:    \> $\texttt{POS(2)} \ge 0$ for \texttt{LCYLX=T}, \texttt{LCYLY=T}, or
                \texttt{LCYLZ=T} \\ 
             \> otherwise: arbitrary
\end{keytab}
\end{keydescription}

\begin{keydescription}{\texttt{POSFIX} --- Fixed target position}
%
  This parameter specifies the position which remains fixed during
  target relaxation in a 1-D dynamic simulation (\texttt{NDIM=1},
  \texttt{LDYN=T}) or during the initial relaxation when the geometry is read 
  from a \texttt{CELL} file (\texttt{NDIM=1}, \texttt{USECELL=T}).  If 
  \texttt{POSFIX} is unspecified, it is assumed in the center of the target.
  The center position is updated during a dynamic simulation. \texttt{POSFIX}
  also falls back to the target center if it happens to fall outside the target
  in the course of a dynamic simulation.  The units of \texttt{POSFIX} are \AA.
  \begin{keytab}
    Type:    \> real \\
   Default:  \> center of simulation domain \\
   Range:    \> arbitrary
\end{keytab}
\end{keydescription}

\begin{keydescription}{\texttt{POSIF} --- Positions of the interfaces}
%
  This parameter specifies the positions of the interfaces between
  neighboring regions in case of a 1-D simulation (\texttt{NDIM=1} on
  the \texttt{\&SETUP} record). The first value of \texttt{POSIF}
  defines the front surface of the target, the last value the back
  surface.  A maximum of 20 positions may be specified on one record.
  Specifications of \texttt{POSIF} on different records are joined together.  In
  this way, larger arrays of positions may be split to several records.  The
  units of \texttt{POSIF} are \AA. 
  \begin{keytab}
    Type:    \> simple array (\texttt{NR+1}) of real \\
    Default: \> 0., 10$^{30}$ if only one region, \\
             \> required otherwise \\
    Range:   \> values must be increasing \\
             \> $\texttt{POSIF(0)} \ge 0$ for \texttt{LSPHER=T}, 
                \texttt{LCYLX=T}, \texttt{LCYLY=T}, or \texttt{LCYLZ=T}
  \end{keytab}
\end{keydescription}

\begin{keydescription}{\texttt{SDFFILE} --- SDF file}
%
  This parameter specifies the name of a signed distance function file. 
  An SDF file must be specified either for all regions or for none. For
  2-D simulations (\texttt{NDIM=2}), the geometry must be specified
  either by SDF files or by \texttt{POS} and \texttt{POINTS}.
  \begin{keytab}
    Type:    \> index variable array (\texttt{REGION}) of character
                strings  \\
    Default: \> -- \\
    Range:   \> any name ($\le$ 80 characters) of a file containing 
                SDF data.
  \end{keytab}
\end{keydescription}

\begin{keydescription}{\texttt{WBOX} --- 1-D box width}
%
  This parameter specifies the cell size of the cell grid used to store
  densities in a dynamic simulation (\texttt{LDYN=T}), if the geometry is 
  defined by regions (\texttt{USECELL=F}). 
  A finer mesh means a finer resolution of the geometry.
  The value of \texttt{WBOX} might be adjusted by IMSIL if the maximum 
  number of boxes \texttt{NBOX} were exceeded, or to fit an integer 
  number of intervals between the boundaries. 
  The units are \AA. 
  \texttt{WBOX} is ignored if the geometry is read from a cell file 
  (\texttt{USECELL}) or if $\texttt{NDIM} \ne 1$.
  \begin{keytab}
    Type:    \> real \\
    Default: \> \texttt{WBOX} of \texttt{\&OUTPUT} for \texttt{LDYN=T} if 
                defined \\
             \> 1.0 otherwise \\
    Range:   \> $> 0$
  \end{keytab}
\end{keydescription}

\begin{keydescription}{\texttt{WBOX2} --- 2-D box width}
%
  This parameter specifies the cell size of the cell grid used to store
  densities in a dynamic simulation (\texttt{LDYN=T}), if the geometry is 
  defined by regions. 
  It is also used to specify the box width of the internal mesh for both
  static (\texttt{LDYN=F}) and dynamic (\texttt{LDYN=T}) simulations if the
  geometry is defined by regions.
  The internal mesh is used to store the signed distance function that
  internally represents the geometry.
  Therefore, a finer mesh means a finer resolution of the geometry.
  The first value of \texttt{WBOX2} corresponds to the cell width in $x$ 
  direction, the second value to the cell width in $z$ 
  direction. 
  The values of \texttt{WBOX2} might be adjusted by IMSIL, if the maximum 
  number of boxes \texttt{NBOX2} were exceeded, or to fit an integer 
  number of intervals between the boundaries. 
  The units are \AA. 
  \texttt{WBOX2} is ignored if the geometry is read from an SDF file 
  (\texttt{SDFFILE}) or from a cell file (\texttt{USECELL}), or for 
  $\texttt{NDIM} \ne 2$.
  \begin{keytab}
    Type:    \> simple array (2) of real \\
    Default: \> \texttt{WBOX2} of \texttt{\&OUTPUT} for \texttt{LDYN=T} if
                defined \\
             \> 1.0, \texttt{WBOX2(1)} otherwise \\
    Range:   \> $> 0$
  \end{keytab}
\end{keydescription}

\begin{keydescription}{\texttt{WROUGH} --- Thickness of roughness layer}
%
  This parameter specifies the thickness of the surface roughness layer.
  Within the roughness layer, the density is assumed to increase linearly
  from 0 (at the surface) to its nominal value (at the other boundary of the
  roughness layer). Note that specifying a $\texttt{WROUGH} > 0$ effectively
  removes material from the target and moves the ``average'' surface
  $\texttt{WROUGH}/2$ into the target. The units are \AA. 
  \begin{keytab}
    Type:    \> real \\
    Default: \> 0 \\
    Range:   \> $\ge 0$
  \end{keytab}
\end{keydescription}

\begin{keydescription}{\texttt{XLIM} --- x boundaries of geometry domain}
%
  This parameter specifies the minimum and maximum $x$ coordinates of the
  geometry domain. It can be used to add vacuum layers left and right to
  the geometry domain if the geometry is defined by polygons 
  (\texttt{POS}, \texttt{POINTS}). Its intended 
  purpose is to add slits to periodic structures. It can also be used to
  increase the vacuum layer on top of a free boundary (which generally is
  useless). \texttt{XLIM} is
  ignored if the geometry is read from an SDF or cell file. In this case, 
  \texttt{XLIM} is set to the limits of $x$ values on the SDF or cell file. 
  If a vacuum layer is desired with an SDF or cell file, the SDF or cell 
  data must be adapted, appropriately. \texttt{XLIM} is ignored unless
  $\texttt{NDIM}=2$. The units are \AA.
  \begin{keytab}
    Type:    \> simple array (2) of real \\
    Default: \> $x$ limits of geometry, except at free boundaries where
                a thin vacuum \\
             \> layer is added  \\
    Range:   \> $\texttt{XLIM(1)} \le \min(\texttt{POS(1)}) <  
                \max(\texttt{POS(1)}) \le \texttt{XLIM(2)}$ 
  \end{keytab}
\end{keydescription}

\begin{keydescription}{\texttt{XSCALE} --- Scaling factor in x direction}
%
  This parameter specifies a scaling factor for the $x$ coordinates of the
  polygon points as specified by \texttt{POS}. All values \texttt{POS(1)} will
  be multiplied by \texttt{XSCALE}. Note that \texttt{WBOX2(1)} will not be 
  scaled. 
  \begin{keytab}
    Type:    \> real \\
    Default: \> 1 \\
    Range:   \> arbitrary
  \end{keytab}
\end{keydescription}

\begin{keydescription}{\texttt{ZLIM} --- z boundaries of geometry domain}
%
  This parameter specifies the minimum and maximum $z$ coordinates of the
  geometry domain. It can be used to add vacuum layers left and right to
  the geometry domain if the geometry is defined by interface positions
  (\texttt{POSIF}) or polygons (\texttt{POS}, \texttt{POINTS}). Its intended 
  purpose is to add slits to periodic structures. It can also be used to
  increase the vacuum layer on top of a free boundary (which generally is
  useless). \texttt{ZLIM} is
  ignored if the geometry is read from an SDF or cell file. In this case, 
  \texttt{ZLIM} is set to the limits of $z$ values on the SDF or cell file. 
  If a vacuum layer is desired with an SDF or cell file, the SDF or cell 
  data must be adapted appropriately. The units are \AA.
  \begin{keytab}
    Type:    \> simple array (2) of real \\
    Default: \> $z$ limits of geometry, except at free boundaries where
                a thin vacuum \\
             \> layer is added  \\
    Range:   \> $\texttt{ZLIM(1)} \le \min(\texttt{POSIF}) <  
                \max(\texttt{POSIF}) \le \texttt{ZLIM(2)}$, \\
             \> $\texttt{ZLIM(1)} \le \min(\texttt{POS(2)}) <  
                \max(\texttt{POS(2)}) \le \texttt{ZLIM(2)}$ \\
             \> $\texttt{ZLIM(1)} \ge 0$ for \texttt{LCYLX=T}, 
                \texttt{LCYLY=T}, \texttt{LCYLZ=T}, or \texttt{LSPHER=T}
  \end{keytab}
\end{keydescription}

\begin{keydescription}{\texttt{ZSCALE} --- Scaling factor in z direction}
%
  This parameter specifies a scaling factor for the $z$ coordinates of the
  polygon points as specified by \texttt{POS}. All values \texttt{POS(2)}
  will be multiplied by \texttt{ZSCALE}. Note that \texttt{WBOX2(2)} will not be
  scaled.
  \begin{keytab}
    Type:    \> real \\
    Default: \> 1 \\
    Range:   \> arbitrary
  \end{keytab}
\end{keydescription}

